\documentclass[a4paper,12pt,oneside]{article}

\usepackage[T2A]{fontenc}
\usepackage[utf8]{inputenc}
\usepackage[russian]{babel}

\usepackage{microtype}
\usepackage{graphicx}

\usepackage{amsmath,amssymb,amsthm,mathtools}

\usepackage{geometry}
\geometry{top=25mm,left=25mm,right=20mm,bottom=25mm}

\usepackage{enumitem}
\setlist[enumerate]{leftmargin=*, itemsep=0.55\baselineskip, topsep=0.35\baselineskip}
\setlist[itemize]{leftmargin=*, itemsep=0.25\baselineskip, topsep=0.25\baselineskip}

\setlength{\parindent}{1.25em}
\setlength{\parskip}{0pt}

\newcommand{\Task}[1]{\section*{\normalfont\bfseries Задание #1}}
\newcommand{\NOT}[1]{\overline{#1}}
\newcommand{\NOR}{\downarrow}
\newcommand{\NAND}{\mid}

\begin{document}

\begin{center}
    {\Large\bfseries Домашняя работа №3. Булевы функции}\\[2mm]
\end{center}
\vspace{2mm}
\hrule
\vspace{4mm}

\Task{1}
Получить СДНФ и СКНФ для функции. Представить функцию в виде вектора минтермов и вектора макстермов.
\begin{enumerate}[label=\arabic*)]
\item $f_1 = x \leftrightarrow \bigl(y \NOR \NOT{z}\bigr)$.
\item $f_2 = (10101001)$.
\item $f_3(A,B,C,D) = A + BC + D$.
\item $f_4 = f_{47541}^{(4)}$.
\end{enumerate}

\Task{2}
Разложить функции по всем переменным.\\
В ответ записать вектор минтермов для $f_1$ и вектор макстермов для $f_2$.
\[
f_1(A,B,C,D)=A\NOT{C}+A\NOT{B}+\NOT{A}C+D,
\]
\[
f_2(A,B,C,D)=(A+\NOT{C})\,(A+\NOT{B})\,(\NOT{A}+C)\,D.
\]

\Task{3}
Функция задана вектором минтермов. Методом Квайна найдите сокращённую ДНФ. Укажите число простых импликант, из которых состоит сокращённая форма, и число вхождений переменных. Постройте карту Карно и по ней найдите тупиковые и минимальные ДНФ и КНФ.
\[
f(A,B,C,D)=(0,3,5,6,7,9,10,13,14,15).
\]

\Task{4}
Найти СДНФ, СКНФ, сокращённую, тупиковые и минимальные ДНФ и КНФ при помощи карты Карно.
\[
\includegraphics[width=0.5\linewidth]{HW3.map.png}
\]

\Task{5}
Перечислить все существенные и фиктивные переменные у функции:
\[
f=(1011010110110101).
\]

\Task{6}
Докажите, что любую функцию, кроме тождественной $1$, можно представить в виде СКНФ.

\Task{7}
Построить полином Жегалкина для следующих формул. \\
\underline{Для каждой формулы выбрать один метод построения полинома Жегалкина, но должны} \\
\underline{быть представлены хотя бы три различных метода построения!}
\begin{enumerate}[label=\alph*)]
\item $x \rightarrow (\neg y \land z)$;
\item $\neg (x \lor y) \NOR z$;
\item $(x \to y)\ \NAND\ \neg(y \land \neg z)$;
\item \(\bigl((x \to y)\ \NAND\ \neg x\bigr)\land z\).
\end{enumerate}

\Task{8}
Проверить, является ли данная система функций полной. В случае, если да, выяснить, является ли она безызбыточной.
\begin{enumerate}[label=\alph*)]
\item $\{\,x \leftrightarrow y,\ x \lor y\,\}$;
\item $\{\,x \lor y,\ (x \leftrightarrow y)z,\ x \NOR (yz)\,\}$;
\item $\{\,\neg x,\ x \lor \neg y\,\neg z\,\}$;
\item $\{\,x \leftrightarrow y,\ \neg x,\ \neg x \leftrightarrow y\,\}$;
\item $\{\,x \NOR \neg y,\ \neg x \NAND y,\ x \to y\,\}$.
\end{enumerate}

\Task{9}
Проверить следующее утверждение: $X \subseteq Y$, где $X$, $Y$ – классы функций, сохраняющие нуль, сохраняющие единицу и монотонные.

\end{document}
